\documentclass[14pt]{beamer}
\usetheme{Berkeley}
\usepackage[utf8]{inputenc}
\usepackage[ngerman]{babel}
\usepackage[T1]{fontenc}
\usepackage{amsmath}
\usepackage{amsfonts}
\usepackage{amssymb}
\usepackage{graphicx}
%\usepackage{svg}
%\usepackage{pdfplots}
%\usepackage{pdfplotstable}
\graphicspath{{images/}}
\author{AUTHOR}
\title{TITLE}
\usepackage{mathptmx}
%\setbeamercovered{transparent} 
%\setbeamertemplate{navigation symbols}{} 
\logo{\includegraphics[scale=0.4]{logo.png}} 
\institute{Herder Gymnasium Berlin} 
\date{} 
%\subject{test} 
\begin{document}

\begin{frame}{Normalverteilung}
\titlepage
\end{frame}

\begin{frame}{Inhalt}
\tableofcontents
\end{frame}

\section{Dichtefunktionen}
\begin{frame}{Wozu eine Dichtefunktion?}
\textbf{Bisherige Berechnung}:

{\small $$ P(X \in\mathbb{A}) \, = \, \sum_{x\in A} \rho(x) $$}
\begin{itemize}

\item eignet sich nur für diskrete Wahrscheinlichkeitsverteilungen:

\begin{itemize}
\only<1>{
\item Binomial
\item Hypergeometrisch
\item Poisson
}
\only<2>{
\item[$\Rightarrow$] keine Aussage möglich für stetige Verteilungen, da $P(X \, = \, x) = 0 \, \forall X \in \mathbb{A}$
}
\end{itemize}
\end{itemize}

\end{frame}

\subsection{Definition}
\begin{frame}{Definition}
\textbf{Wahrscheinlichkeit als Flächeninhalt:}
$\varphi \colon \mathbb{R} \rightarrow [0,\infty)$ mit
$$P(X \in \mathbb{A})=\int_{\mathbb{A}} \varphi(x)\,{\mathrm d}x$$
ist die \textbf{Wahrscheinlichkeitsdichte} der Verteilung von $X$.
\end{frame}

\subsection{Eigenschaften}
\begin{frame}{Eigenschaften}
Eigenschaften:
\begin{itemize}
\item $\forall x \in \mathbb{R}:\, \varphi(x)\geq 0$
\item $\int_{-\infty}^\infty f(x)\,{\mathrm d}x = 1$
\end{itemize}
Eine Funktion $f$ ist genau dann eine Dichtefunktion, wenn diese Eigenschaften erfüllt sind.
\end{frame}

\subsection{Nutzung}
\begin{frame}{Rechnen mithilfe der Dichtefunktion}
Wahrscheinlichkeit für ein Intervall $[a,b]$:
$$P(X\in [a,b]) = \int_a^b f(x)\, {\rm d}x = \Phi(b)-\Phi(a).$$
$F$ ist die sogenannte \textbf{kumultative Verteilungsfunktion} mit:
$$\Phi(x) \, = \, P(X \le x)$$
\end{frame}

\section{Approximation der Binomialverteilung}
\begin{frame}{Was ist $\varphi(x)$?}
\textbf{Ziel: Annäherung an die Binomialverteilung:}
\pause
\begin{itemize}
\item Verschieben des Erwartungswertes, sodass $\mu=0$:\\
$k \to k-\mu$
\item Vereinheitlichen der Breite:\\
$k-\mu \to \frac{k-\mu}{\sigma}$
\pause
\item \glqq Wiederherstellen\grqq\ des Flächeninhalts:\\
$B(n;p;k) \to \sigma \cdot B(n;p;k)$
\end{itemize}

\end{frame}

\begin{frame}{Zwischenstand}
$$
\sigma \cdot B(n;p;k) =: \varphi\ \left(\ \frac{k-\mu}{\sigma}\right)
$$
$$
\Leftrightarrow
\begin{mathbf}
B(n;p;k) =  \frac{1}{\sigma} \cdot  \varphi\ \left(\ \frac{k-\mu}{\sigma}\right)
\end{mathbf}
$$
\end{frame}


\begin{frame}{Zusammenfassung}
$$
P(X\leq\ c) = \frac{1}{\sigma} \cdot \int^c_{-\infty} \varphi\ \left(\frac{t-\mu}{\sigma}\right) {\mathrm d}t
$$
mit Substitution $g(t):=\frac{t-\mu}{\sigma} \Rightarrow g'(t)=\frac{1}{\sigma}$:
$$
= \int^{\frac{c-\mu}{\sigma}}_{-\infty} \varphi(x) {\mathrm d}x =: \phi\ \left(\frac{c-\mu}{\sigma}\right) = \phi\ \left( \frac{c-\mu}{\sqrt{npq}}\right)
$$
\end{frame}

\section{Herleitung: Die Normalverteilung}
\begin{frame}{\textrm{\textsc{Gau}ß}}
Nach \textrm{\textsc{Gau}ß} ist eine dazu passende Funktion $\varphi$:
$$
\varphi(x)=b\cdot e^{-a\cdot\ x^2}
$$
Durch unsere Vereinheitlichung der Breite ist das $1\cdot \sigma$-Intervall $[-1;1]$.
\end{frame}


\begin{frame}{Was ist $a$?}
Betrachten des Graphen der Binomialverteilung zeigt uns, an $x=\pm 1$ sind Wendestellen.\\
Für $x=\pm 1$:\\
2. Ableitung: $2abe^{-a} \cdot (-1+2 a)$\\
3. Ableitung: $-4 a^{2} b e^{-a} (\pm1) (-3+2 a)$\\

Da $\varphi''(x)=0$ und $\varphi'''(x)\neq0$ muss:\\
$a=\pause\frac{1}{2}$
\end{frame}

%\begin{frame}{Was ist $b$?}
%Der Flächeninhalt des Graphen muss $1$ sein:
%$$
%1=\int^{\infty}_{-\infty} b \cdot e^{-\frac{1}{2}x^2} {\mathrm d}x
%$$
%$$
%\Leftrightarrow \frac{1}{b}=\int^{\infty}_{-\infty} e^{-\frac{1}{2}x^2} {\mathrm d}x
%$$
%\end{frame}

\begin{frame}{\textrm{\textsc{Gau}ß}sche Normalverteilung}
Über mehrdimensionale Analysis können wir $b:=\frac{1}{\sqrt{2\cdot \pi}}$ herleiten, damit:
$$
1= \Phi =  \int^{\infty}_{-\infty} \varphi(x) {\mathrm d}x = \int^{\infty}_{-\infty} \frac{1}{\sqrt{2\cdot \pi}} \cdot e^{-\frac{1}{2}x^2} {\mathrm d}x
$$
$\varphi(x)$: \textrm{\textsc{Gau}ß}sche \textbf{Dichtefunktion}\\
$\Phi(x)$:  \textrm{\textsc{Gau}ß}sche \textbf{Summenfunktion}
\end{frame}



\section{Was sagt der Grenzwertsatz aus?} % (Bedeutung)
\subsection{Beweis} %  od. Plausibilitätsprüfung
\section{Gauß} %Hinführung auf Gauß'sche Formel (e^(-1/2x^2)), so wie es Herr Steinkrauß erläutert hat (oder auch anders, aber in diesem Detailgrad)
\section{Aufgaben zur lokalen Näherung}
\end{document}




Was ich dann (nachfolgend) mache:
- Was sagt das Integral aus?
- Wie bilde ich überhaupt das Integral?
- Wie gucke ich die Werte nach? (Tabellen)
+ ein paar Aufgaben 